\section{Khảo sát hiện trạng}

\subsection{Đặt vấn đề}

Máy chủ trò chuyện (Chat Server) đóng vai trò trung tâm trong hệ thống truyền thông thời gian thực (Real-time Communication), đảm nhiệm chức năng định tuyến thông điệp song công (Full-duplex) giữa các máy khách (Client).

Trong điều kiện xử lý tải cao, mô hình cấp phát luồng theo kết nối (Thread-per-client) không đáp ứng hiệu quả bài toán đa kết nối đồng thời (C10K), dẫn đến thắt cổ chai I/O (I/O Bottleneck) và suy giảm hiệu suất sử dụng tài nguyên. Bên cạnh đó, việc sử dụng giao thức văn bản (Text-based Protocol) phát sinh chi phí phân tích cú pháp (Parsing overhead) và gia tăng dung lượng truyền tải. Để khắc phục các vấn đề này, giải pháp kết hợp mô hình I/O đa kênh (I/O Multiplexing) và giao thức nhị phân (Binary Protocol) được áp dụng nhằm tối ưu hóa băng thông và giảm thiểu độ trễ.

Bảo vệ an toàn dữ liệu truyền tải là thành phần thiết yếu của hệ thống. Sự thiếu hụt cơ chế kiểm soát truy cập và bảo vệ kênh truyền làm tăng nguy cơ từ các hình thức tấn công nghe lén (Eavesdropping), tấn công xen giữa (MitM) và từ chối dịch vụ (DoS/Spam). Do đó, hệ thống yêu cầu tích hợp các thành phần bảo mật bao gồm cơ chế mã hóa dữ liệu, hệ thống phát hiện xâm nhập (IDS) và giới hạn lưu lượng (Rate Limiter).

Dựa trên cơ sở đó, nghiên cứu này tập trung phát triển hệ thống máy chủ trò chuyện an toàn (Secure Chat Server) hiệu năng cao, hướng tới tối ưu hóa I/O mạng, quá trình định tuyến thông điệp và triển khai cơ chế bảo mật toàn diện.

\subsection{Mục tiêu và phạm vi của đề tài}
\begin{enumerate}
\item {Mục tiêu của đề tài} \\
Đề tài nghiên cứu và phát triển hệ thống máy chủ truyền thông (Secure Chat Server) theo mô hình Client-Server với các mục tiêu cụ thể:
\begin{itemize}[label={-}]
    \item Tối ưu hóa hiệu năng: Triển khai cơ chế I/O Multiplexing (epoll) kết hợp Thread Pool để giải quyết bài toán C10K, cải thiện thông lượng xử lý I/O và giảm độ trễ hệ thống.
    \item Tối ưu hóa giao thức: Áp dụng giao thức nhị phân tùy chỉnh (Custom Binary Protocol) nhằm loại bỏ chi phí phân tích cú pháp và tối thiểu hóa băng thông mạng.
    \item Đảm bảo an toàn thông tin: Tích hợp cơ chế mã hóa lai (RSA-2048, AES-256-GCM), hệ thống phát hiện xâm nhập (IDS), thuật toán giới hạn lưu lượng (Token-Bucket Rate Limiter) và ghi log kiểm toán (Audit Logger).
    \item Đảm bảo tính nhất quán dữ liệu: Tích hợp cơ sở dữ liệu SQLite cùng cơ chế Write-Ahead Logging (WAL) để tối ưu hóa hiệu suất truy vấn trong môi trường đa luồng.
\end{itemize}

\item {Phạm vi nghiên cứu và triển khai} \\
Phạm vi nghiên cứu và giới hạn của đề tài bao gồm:
\begin{itemize}[label={-}]
    \item Kiến trúc hệ thống: Lõi máy chủ (Server-side) được lập trình bằng ngôn ngữ C++. Ứng dụng máy khách (Client-side) được triển khai dưới dạng giao diện dòng lệnh (CLI) phục vụ mục đích kiểm thử (giới hạn không bao gồm giao diện người dùng đồ họa hoặc ứng dụng di động).
    \item Chức năng nghiệp vụ: Hỗ trợ tính năng nhắn tin cá nhân (Direct Message), trò chuyện nhóm (Rooms), truyền tải tệp tin (File Transfer) và tích hợp các tập lệnh quản trị (Admin Commands).
    \item Đánh giá hệ thống: Thực hiện đo lường hiệu năng (Benchmark) dưới tải cao và mô phỏng các kịch bản tấn công (DoS, quét cổng, giả mạo gói tin) nhằm đánh giá năng lực phòng thủ.
\end{itemize}
\end{enumerate}
